\chapter*{Conclusion}

\addcontentsline{toc}{chapter}{Conclusion}

Shelfwise implements an inventory and stock-warning workflow for one supermarket. The completed application connects a current product balance to accepted stock movements and connects shortage interpretation to persisted warning episodes. A responsive Vue interface gives administrators, managers and clerks access to different operations, while a Spring Boot server and MySQL database remain responsible for authorization and durable change. Redis is an optional page cache rather than a second source of business truth.


The first engineering question concerned bounded and explainable quantity changes. The design answers it with a locked product read, server-derived arithmetic, mandatory actor and reason, an actor-key uniqueness constraint and a coordinated transaction. The integration evidence includes receipt, dispatch, signed correction, absolute counts, rejection without history, exact replay and selected concurrent operations. Seven units cannot be oversold by the twelve competing unit-dispatch attempts in the implemented test. This is a concrete supported case, not a universal workload proof.


The second question concerned the interpretation of shortage and recovery. The system distinguishes OUT, LOW and a time-limited RESTOCKED episode while storing acknowledgement separately. A review records the first reviewer without changing quantity. The tests support zero and equality boundaries, healthy replenishment, renewed shortage, archival and concurrent acknowledgement. Complete time-based expiry and cached visibility at that boundary remain useful additional cases. The distinction between review and replenishment is nevertheless preserved in the implemented model and user-facing controls.


The third question concerned the relationship between interface authority and server authority. Fixed role bundles and endpoint rules separate catalogue administration, stock posting and managerial review. Browser navigation reduces inappropriate paths, while integration checks exercise server denial. Account enabled state and current role are refreshed on subsequent requests. The resulting boundary does not depend on a hidden button being the only obstacle to an unauthorized write. Further deployment work must still address rate limiting, credential recovery, MFA, session scaling and a wider security assessment.


The fourth question concerned evidence for performance and recovery. The preserved warning-query benchmark records a pooled median of 5.33 ms with Redis enabled and 6.56 ms with bypass across six hundred sequential samples per path. Its distributions overlap, and the result includes the revision and authentication work retained in both paths. A recorded Redis interruption returns equivalent warning data, and an ordinary application restart preserves selected stock and movement rows. These observations support the selected local mechanisms without establishing capacity, disaster recovery or a production availability target.


The main outcome is therefore an integrated and inspectable engineering artifact. Requirements, diagrams, implementation excerpts, captured interfaces and executable evidence refer to the same domain concepts. The system makes the difference between a product edit, a movement, a physical count and a review explicit. That explanatory consistency is valuable even before more complex retail functions are introduced.


Future development should first resolve domain limits that affect actual store use. Fractional quantities need exact unit and rounding policies. Batch and expiry tracking need additional stock identities. A point-of-sale integration needs external event identity, retry behavior and reconciliation. Purchasing requires supplier terms and order lifecycle data. These additions should extend the ledger and acceptance catalog before they add dashboard claims.


Operational evaluation should then examine staff workflows, wider concurrent workloads, cache-expiry boundaries, cross-product key races, database deadlocks and restore drills. Frontend decomposition and service refactoring can improve maintenance, provided that the effective transaction and permission boundaries remain intact. The present application supplies a tested baseline for that work, with its limits stated alongside its supported results.
